\chapter{METODOLOGIA CIENTÍFICA}

% A metodologia em um projeto de TCC descreve o plano detalhado de como a pesquisa será conduzida, incluindo os métodos, técnicas, procedimentos e abordagens que serão utilizados para coletar e analisar dados. Ela é uma parte crucial do projeto, pois fornece uma estrutura sistemática para garantir que os objetivos da pesquisa sejam alcançados de forma eficaz e confiável.

\section{Abordagem da pesquisa}

Esta pesquisa possui abordagem predominantemente quantitativa, uma vez que serão realizadas medições e comparações de desempenho relacionadas às estruturas de aceleração utilizadas em algoritmos hibridos de rasterizacao e \textit{ray tracing}. Serão analisadas métricas como tempo de construção, tempo de travessia, uso de memória e desempenho geral da renderização em diferentes abordagens de \gls{bvh}.

A partir dos resultados obtidos, também será realizada uma análise qualitativa, buscando interpretar os impactos das diferentes técnicas na qualidade da hierarquia gerada, na complexidade de implementação e na adequação de cada abordagem para diferentes tipos de cena, aplicações e dispositivos. Dessa forma, a pesquisa combina a análise objetiva de dados experimentais com uma discussão descritiva dos resultados observados.


% A metodologia inclui várias seções importantes, tais como:
% \begin{enumerate}
%     \item Abordagem da pesquisa: Descreve se a pesquisa será qualitativa, quantitativa ou mista. Isso envolve a escolha entre abordagens como estudos de caso, experimentos, levantamentos, entre outros.

%     \item Instrumentos de coleta de dados: Indica quais instrumentos serão utilizados para coletar dados, como questionários, entrevistas, observações, análise de documentos, entre outros. Essa seção também descreve como esses instrumentos serão desenvolvidos e validados, se necessário.

%     \item Procedimentos de coleta de dados: Detalha como os dados serão coletados, incluindo informações sobre o local, o período de tempo, os participantes (se houver), os critérios de seleção, entre outros aspectos relevantes.

%     \item Procedimentos de análise de dados: Descreve como os dados serão processados, analisados e interpretados. Isso pode incluir técnicas estatísticas, análise de conteúdo, codificação de dados, entre outras abordagens.

%     \item Considerações éticas: Discute as questões éticas relacionadas à pesquisa, como consentimento informado dos participantes, confidencialidade, anonimato, entre outros princípios éticos a serem seguidos durante a condução da pesquisa.
% \end{enumerate}

% Para trabalhos voltados ao \textbf{desenvolvimento de sistemas} é esperado:
% \begin{enumerate}
%     \item Levantamento de requisitos
%     \item Definição de tecnologias/infraestrutura
%     \item Modelo de Banco
%     \item Interfaces (prototipação)
% \end{enumerate}



\section{Orçamento}
Se achar necessário, pode citar algo sobre uso ou não de recursos orçamentários